Subbab~\ref{sec:statistik_terintegrasi} telah menyatukan hasil uji statistik
dari seluruh metrik kontinu ke dalam satu pandangan menyeluruh. Pola yang
muncul di sana, peningkatan keselamatan yang konsisten, \textit{trade-off}
efisiensi yang terukur, dan perubahan kualitas gerak yang tidak seragam,
perlu dipahami bukan hanya sebagai angka, tetapi sebagai konsekuensi dari
rancangan metode yang telah diuraikan pada BAB III. Subbab ini membahas
mengapa pola-pola tersebut muncul dan apa maknanya terhadap tujuan
penelitian: apakah integrasi \textit{virtual blind-spot scan}, fusi
\textit{traversability} LiDAR--termal, dan adaptasi fuzzy Mamdani benar-benar
memberikan manfaat terhadap keselamatan navigasi tanpa mengorbankan performa
secara berlebihan.

\subsection{Menjawab Pertanyaan Penelitian}

\textbf{Apakah metode usulan mempertahankan keberhasilan navigasi?}
Berdasarkan Subbab~\ref{sec:keberhasilan_tabrakan}, metode usulan mencapai
\textit{success rate} dan \textit{collision-free rate} 100\% pada seluruh
empat skenario fuzzy (S5--S8). Metode \textit{baseline} juga mencapai 100\%
pada tiga dari empat skenarionya (S1--S3), dengan satu pengecualian pada S4
(90\%, satu dari sepuluh \textit{trial} mengalami tabrakan dengan
\texttt{obstacle\_human\_1} meski robot tetap mencapai \textit{goal}). Data
ini tidak mendukung narasi bahwa \textit{baseline} "gagal", baseline tetap
berhasil pada sembilan dari sepuluh kombinasi skenario-konfigurasi, namun
metode usulan berhasil menutup satu-satunya celah kegagalan yang muncul pada
kondisi paling menantang, yaitu kombinasi \textit{obstacle} dinamis dan
\textit{blind-spot}.

\textbf{Apakah metode usulan meningkatkan keselamatan?} Jawabannya, dengan
catatan, ya. Pada ketiga pasangan yang dapat diuji secara statistik
(Subbab~\ref{sec:analisis_keselamatan}), minimum clearance meningkat secara
signifikan dengan \textit{effect size} besar, dan besarnya peningkatan
absolut justru membesar seiring kompleksitas skenario: dari +0,007~m pada
S2/S6 (lingkungan dengan satu \textit{obstacle} statis biasa), menjadi
+0,121~m pada S3/S7 dan +0,303~m pada S4/S8, dua kondisi yang melibatkan
kucing pada area \textit{blind-spot}. Pola ini sejalan dengan rancangan
metode usulan yang memang ditujukan untuk kondisi tersebut, bukan untuk
lingkungan sederhana.

\textbf{Apakah peningkatan keselamatan dibayar dengan penurunan efisiensi?}
Ya, dan \textit{trade-off} ini konsisten di seluruh pasangan
(Subbab~\ref{sec:tradeoff}). Waktu tempuh meningkat 15,5\% hingga 63,7\%,
dengan kenaikan terbesar pada S3/S7 dan S4/S8, skenario yang sama yang
menunjukkan peningkatan clearance paling besar. Panjang lintasan relatif
tidak banyak berubah ($-0,7\%$ hingga $+5,9\%$), sehingga \textit{trade-off}
utama terletak pada waktu dan kecepatan, bukan pada rute geometris yang
ditempuh. Besarnya biaya ini perlu dibaca berdampingan dengan manfaat
keselamatan yang diperoleh: pada S3/S7, kenaikan waktu 63,7\% disertai
kenaikan clearance 198,4\%, sementara pada S2/S6, kenaikan waktu 15,5\%
hanya disertai kenaikan clearance 1,6\%. \textit{Trade-off} pada pasangan
kedua ini secara proporsional kurang menguntungkan dibandingkan pasangan
pertama, meskipun kedua kombinasi signifikan secara statistik.

\textbf{Apakah kualitas gerak tetap dapat diterima?} Hasil pada
Subbab~\ref{sec:kehalusan_gerak} menunjukkan bahwa peningkatan keselamatan
tidak menyebabkan perilaku navigasi yang ekstrem. Kecepatan angular rata-rata
metode usulan justru lebih rendah pada S1/S5 dan S2/S6, sedikit lebih tinggi
pada S3/S7, dan tidak berbeda signifikan pada S4/S8. Tidak satu pun trajektori
yang diperiksa secara visual menunjukkan pola zig-zag tajam pada konfigurasi
fuzzy, satu-satunya pola zig-zag tajam yang teramati justru muncul pada satu
\textit{trial baseline} di S2. Variasi yang lebih besar pada S7 dibandingkan
S3 lebih tepat digambarkan sebagai perbedaan konsistensi antar-\textit{trial},
bukan ketidakstabilan gerak.

\subsection{Evaluasi Kontribusi Setiap Komponen}

Penelitian ini tidak melakukan \textit{ablation study}, sehingga kontribusi
ketiga komponen metode usulan tidak dapat dipisahkan secara independen dari
data eksperimen semata. Penting dicatat pula bahwa pada tingkat implementasi,
ketiga komponen ini juga tidak benar-benar independen satu sama lain:
penelusuran terhadap \texttt{APFPlanner} (Subbab~\ref{sec:arsitektur})
menunjukkan bahwa mekanisme \textit{cat keep-out} dan fusi
\textit{traversability} via operator minimum hanya aktif bersamaan dengan
adaptasi fuzzy, dikendalikan oleh satu parameter \texttt{use\_fuzzy} yang
sama. Dengan keterbatasan ini sebagai konteks, evaluasi berikut bersifat
naratif berdasarkan rancangan, bukan klaim kausal yang teruji.

\textit{Virtual blind-spot scan} (proyeksi blob termal menjadi
\texttt{sensor\_msgs/LaserScan} pada Subbab~\ref{sec:proyeksi_scan})
memungkinkan informasi kucing yang tidak terjangkau LiDAR dimanfaatkan oleh
\textit{costmap} dan \textit{planner} berbasis \texttt{LaserScan}, kompatibel
langsung dengan ekosistem ROS \textit{navigation stack} tanpa memerlukan
modifikasi pada \textit{global planner}. Mekanisme ini secara desain aktif
pada kedua konfigurasi, bukan eksklusif fuzzy. Fakta bahwa \textit{baseline}
(S3) tetap mencapai 100\% keberhasilan dan \textit{collision-free} di sekitar
kucing, meski dengan margin yang jauh lebih tipis (0,061~m) dibandingkan
fuzzy (S7, 0,182~m), menjadi indikasi tidak langsung bahwa jalur informasi
ini memang memberikan manfaat dasar pada kedua konfigurasi, tanpa virtual
scan, LiDAR semata tidak memiliki cara untuk mengetahui keberadaan kucing
sama sekali.

Fusi \textit{traversability} (Subbab~\ref{sec:fusi_trav}) menggabungkan skor
geometris dari LiDAR dengan skor termal melalui operator minimum
($\tau_\text{eff} = \min(\tau_\text{LiDAR}, \hat{\tau}_\text{thermal})$),
menghasilkan representasi kondisi lingkungan yang secara desain selalu
mengambil estimasi yang lebih konservatif dari kedua sumber. Mekanisme ini
hanya digunakan secara bermakna pada konfigurasi fuzzy; \textit{baseline}
hanya menggunakan skor termal mentah pada formula \textit{gain crisp}-nya.
Pola peningkatan clearance yang konsisten pada S3/S7 dan S4/S8 sejalan
dengan desain operator minimum ini, meskipun karena keterikatannya dengan
saklar fuzzy yang sama, kontribusinya tidak dapat dipisahkan dari adaptasi
fuzzy itu sendiri.

Adaptasi fuzzy Mamdani memungkinkan \texttt{speed\_scale} dan
\texttt{rep\_gain} berubah secara adaptif berdasarkan jarak ke \textit{obstacle}
terdekat dan $\tau_\text{eff}$, alih-alih formula \textit{crisp} tetap pada
\textit{baseline}. Efeknya tampak konsisten pada pola kecepatan linear yang
selalu lebih rendah pada konfigurasi fuzzy (Subbab~\ref{sec:tradeoff}) dan
pada peningkatan clearance yang telah dibahas. Namun, mengingat fuzzy juga
menjadi saklar yang mengaktifkan \textit{cat keep-out} dan penggunaan
bermakna $\tau_\text{eff}$, peningkatan yang teramati pada eksperimen ini
sebaiknya dipahami sebagai efek gabungan ketiga komponen yang bekerja
sebagai satu kesatuan, bukan kontribusi adaptasi fuzzy semata.

\subsection{Hubungan dengan Literatur}

Karakteristik APF yang rentan terhadap osilasi di sekitar \textit{obstacle}
pada \textit{local minima} telah dikenal sejak diperkenalkan oleh
Khatib~\cite{Khatib1986} sebagai metode penghindaran reaktif berbasis gaya
tarik-tolak. Penelitian ini tidak menghilangkan kerentanan tersebut secara
fundamental, mekanisme deteksi osilasi pada Persamaan~\ref{eq:osc_detect}
(Subbab~\ref{sec:algoritma_navigasi}) tetap berupa mitigasi reaktif, bukan
perubahan formulasi APF itu sendiri seperti pendekatan \textit{subarea-APF}
yang diusulkan Lv dkk.~\cite{Lv2024}. Temuan pada
Subbab~\ref{sec:kehalusan_gerak}, bahwa kecepatan angular fuzzy justru lebih
tinggi pada S3/S7, konsisten dengan pemahaman bahwa adaptasi fuzzy pada
penelitian ini bekerja sebagai pengatur kecepatan dan penguatan repulsi,
bukan sebagai solusi struktural atas osilasi APF.

Pada sisi fusi sensor, pendekatan menggabungkan LiDAR dan kamera termal
untuk mengatasi keterbatasan deteksi pada kondisi tertentu telah ditunjukkan
pada penelitian fusi LiDAR--termal untuk kendaraan otonom~\cite{Choi2023}
dan untuk lokalisasi pada lingkungan berbahaya~\cite{Schichler2025}. Hasil
penelitian ini sejalan dengan temuan tersebut: kombinasi kedua modalitas
sensor memungkinkan deteksi objek yang tidak terjangkau salah satu sensor
saja. Pada konteks estimasi \textit{traversability} secara khusus, survei
oleh Sevastopoulos dan Konstantopoulos~\cite{Sevastopoulos2022} menyoroti
pentingnya representasi kondisi lingkungan yang konservatif untuk keperluan
navigasi aman, sejalan dengan pemilihan operator minimum pada fusi
\textit{traversability} penelitian ini.

Pada sisi navigasi berbasis fuzzy, peningkatan keselamatan dengan
\textit{trade-off} efisiensi yang masih dapat diterima juga ditemukan pada
penelitian navigasi berbasis \textit{neuro-fuzzy} dan fusi sensor pada
lingkungan padat \textit{obstacle}~\cite{Haider2022}, serta pada pengendali
fuzzy logic ganda untuk navigasi pada lingkungan dinamis yang divalidasi
baik secara simulasi maupun pada robot nyata Pioneer 3-DX~\cite{Benaicha2026}.
Pola umum bahwa adaptasi fuzzy meningkatkan keamanan dengan konsekuensi pada
efisiensi gerak, sebagaimana dilaporkan kedua penelitian tersebut, konsisten
dengan apa yang ditemukan pada penelitian ini, meskipun besaran
\textit{trade-off} yang dilaporkan tidak dapat dibandingkan secara langsung
karena perbedaan platform, skenario uji, dan metrik yang digunakan.

\subsection{Kontribusi Penelitian}

Kontribusi metodologis penelitian ini meliputi integrasi blob termal menjadi
\textit{virtual LaserScan} yang kompatibel langsung dengan \textit{costmap}
ROS tanpa modifikasi \textit{global planner} (Subbab~\ref{sec:proyeksi_scan}),
fusi \textit{traversability} LiDAR--termal melalui operator minimum dengan
\textit{low-pass filter} asimetris yang secara desain konservatif terhadap
keselamatan (Subbab~\ref{sec:fusi_trav}), dan adaptasi APF menggunakan fuzzy
Mamdani yang mengendalikan dua parameter sekaligus (\texttt{speed\_scale} dan
\texttt{rep\_gain}) berdasarkan kombinasi jarak \textit{obstacle} dan
\textit{traversability} efektif (Subbab~\ref{sec:fuzzy_mamdani}).

Kontribusi empiris meliputi peningkatan minimum clearance yang signifikan
dan \textit{effect size} besar pada seluruh pasangan yang dapat diuji,
particularly besar pada skenario \textit{blind-spot} dan dinamis;
keberhasilan menutup satu-satunya celah kegagalan \textit{baseline} pada S4;
serta kuantifikasi \textit{trade-off} efisiensi yang terukur dan dapat
dipertimbangkan berdasarkan konteks aplikasi, bukan diklaim sebagai
keunggulan mutlak.

\subsection{Keterbatasan Penelitian}

Beberapa keterbatasan perlu dicatat secara objektif. Seluruh eksperimen
dilakukan pada simulasi Gazebo dan belum diuji pada robot fisik, sehingga
faktor seperti \textit{noise} sensor nyata, dinamika kontak, dan variasi
pencahayaan terhadap kamera termal belum tercakup. Penelitian ini tidak
melakukan \textit{ablation study}, sehingga kontribusi relatif ketiga
komponen metode usulan, yang juga secara arsitektural saling terikat pada
saklar \texttt{use\_fuzzy} yang sama, tidak dapat dipisahkan. Jumlah
\textit{trial} per skenario ($n=10$) relatif kecil untuk standar statistik
umum, meskipun telah dipertimbangkan melalui uji non-parametrik dan
pelaporan \textit{effect size}. Penelitian ini juga hanya menggunakan satu
konfigurasi parameter fuzzy yang telah di-\textit{tuning} untuk delapan
skenario spesifik, sehingga generalisasi terhadap \textit{rule base} atau
\textit{membership function} yang berbeda, maupun terhadap skenario dengan
kompleksitas lebih tinggi (misalnya \textit{obstacle} dinamis majemuk),
belum dapat dipastikan.

\subsection{Implikasi Penelitian}

Secara akademik, penelitian ini menyediakan dasar empiris mengenai bagaimana
fusi sensor LiDAR--termal dan adaptasi fuzzy dapat saling melengkapi pada
navigasi berbasis APF, khususnya untuk kondisi \textit{blind-spot} yang
relatif sedikit dibahas pada literatur navigasi mobile robot berbasis LiDAR
tunggal. Secara praktis, pendekatan ini berpotensi relevan untuk robot
layanan yang beroperasi pada lingkungan semi-terstruktur dengan kemungkinan
kehadiran objek bersuhu tubuh pada area yang tidak terjangkau LiDAR, seperti
fasilitas publik atau area dalam ruangan dengan lalu lintas manusia maupun
hewan, meskipun penerapan langsung pada robot nyata memerlukan validasi
lebih lanjut di luar lingkup penelitian ini.

\subsection{Arah Pengembangan Selanjutnya}

Beberapa arah pengembangan yang realistis meliputi pengujian pada robot
fisik untuk memvalidasi temuan simulasi ini, pelaksanaan \textit{ablation
study} dengan memisahkan saklar \textit{cat keep-out}, fusi
\textit{traversability}, dan adaptasi kecepatan/penguatan fuzzy secara
independen, eksplorasi \textit{rule base} fuzzy yang adaptif terhadap
skenario alih-alih \textit{tuning} tetap, serta pertimbangan integrasi
pembelajaran mesin untuk estimasi \textit{traversability} yang lebih
generalis dibandingkan operator minimum berbasis aturan tetap saat ini.

\subsection{Menjawab Hipotesis Penelitian}

Penelitian ini belum memiliki rumusan hipotesis tertulis secara formal pada
BAB I. Sebagai acuan, pertanyaan validasi hipotesis utama yang telah
dirumuskan pada Subbab~\ref{sec:interpretasi}, yaitu apakah konfigurasi
fuzzy secara signifikan meningkatkan clearance minimum dibandingkan
\textit{baseline} pada skenario dengan \textit{obstacle} kucing (S3/S7),
didukung oleh data pada penelitian ini ($p=0,0002$; $r=+1,00$). Tiga
pertanyaan pendukung lainnya pada subbab yang sama, mengenai \textit{trade-off}
waktu, kehalusan gerak, dan ketiadaan penurunan performa pada S1/S5, juga
didukung sebagian besar oleh data, dengan satu pengecualian pada kecepatan
angular S4/S8 yang tidak signifikan. Dukungan ini bersifat parsial dan
spesifik konteks, tidak dapat dinyatakan sebagai pembuktian mutlak mengingat
keterbatasan ukuran sampel, lingkup simulasi, dan ketidakhadiran
\textit{ablation study} yang telah diuraikan di atas.

Setelah seluruh hasil dibahas secara terintegrasi pada subbab ini, subbab
berikutnya merangkum temuan utama penelitian sebelum memasuki BAB V,
disajikan pada Subbab 4.9 Ringkasan Temuan BAB IV.